\documentclass[11pt]{article}
\usepackage{icmlko}
\begin{document}

\title{판이 못 보는 손해 --- 용량을 더 주면 무리 안에서만 잃는다}
\author{월드모델 실험실}
\date{노트 473 · 2026년 8월 3일}
\maketitle

\begin{abstract}
lr $0.10$ · \texttt{it} $220$ · K $32$ 은 셋 다 \textbf{판으로만} 정했다
(노트 $456 \cdot 459 \cdot 460$). 노트 470 이 만든 \textbf{무리 안 판}
으로도 같은 값이 최선인지 안 봤길래 쟀다.

\textbf{챔피언 값은 확인됐다} --- 무리 안 판에서 챔피언을 \emph{이기는}
후보는 없다. 셋($\texttt{it}\,120 \cdot \text{lr}\,0.06 \cdot \text{K}\,64$)
이 $+0.0001 \sim +0.0002$ 로 나왔지만 구간이 모두 $0$ 을 물어 \textbf{비긴
것}이다. (등록할 때 부호만 보고 문턱을 안 잡았다 --- 노트 451 규약을
또 어겼다. 여기 적어 둔다.)

\textbf{진짜 신호는 반대쪽에 있었다.}

\begin{center}
\begin{tabular}{lrr}
\hline
후보 $-$ 챔피언 & 판 차 & 무리 안 차 \\
\hline
\texttt{it} $400$ & $+0.0003$ $[-0.0005,+0.0011]$ & $\mathbf{-0.0024}$ $[-0.0035,-0.0013]$ \\
lr $0.16$ & $+0.0006$ $[-0.0011,+0.0023]$ & $\mathbf{-0.0016}$ $[-0.0037,+0.0005]$ \\
\hline
\end{tabular}
\end{center}

\textbf{용량을 더 주면 판은 비기는데 무리 안에서만 잃는다.} \texttt{it}
$400$ 은 구간이 $0$ 을 안 문다.
\end{abstract}

\begin{figure}[t]
\centering
\includegraphics[width=.85\textwidth]{figs/f1.pdf}
\caption{다섯 후보를 두 자 위에. 오른쪽 아래(붉은 칸)가 \emph{판은
비기는데 무리 안에서 잃는} 자리이고, 용량을 늘린 둘이 거기 있다.}
\end{figure}

\section{노트 461 이 예고한 것이 새 자에 잡혔다}

노트 461 은 \texttt{it} 을 \textbf{외우는 용량}이라 불렀다 --- 학습
적합도를 $+0.124$ 사면서 집 밖을 $-0.023$ 판다. 그때 남은 물음은
\emph{그 외우기가 판에는 왜 안 보이나} 였다. \texttt{it} $400$ 은 판에서
$+0.0003$ 으로 멀쩡하다.

\textbf{답이 여기 있다.} \texttt{it} 을 늘리면 모델은 \textbf{무리를
맞히는 쪽}으로 능력을 옮긴다. 판은 도메인 전체 순위라 그 쪽을 보상하고,
무리 안 판은 안 보상한다. 그래서 \textbf{판에서는 비기고 무리 안에서는
잃는다.}

노트 470 이 무리 안 판을 만들면서 적은 문장 --- \emph{``무리 사이만 잘
맞히는 변경도 판을 속일 수 있다''} --- 의 첫 실물이다. 노트 472 는 이
사이클의 채택 셋에서 그런 것을 못 찾았는데, \textbf{채택하지 않은
쪽에서 나왔다.}

\section{그래서 자가 하나 늘어난 값을 한다}

노트 470 이 무리 안 판을 만들 때는 \emph{기록 지표}였다. 이제 그것이
\textbf{실제로 무엇을 거른다}는 것을 안다.

\begin{quote}
\textbf{조항} --- 용량을 늘리는 변경은 \textbf{무리 안 판을 반드시 같이
잰다.} 판이 비기는데 무리 안 판이 내려가면, 그 변경은 \emph{무리를
맞히는 쪽으로 능력을 옮긴 것}이므로 넣지 않는다. (지금 규약의 조건
①$\sim$④ 는 이것을 못 잡는다 --- 판이 안 움직이면 조건 ①이 침묵하고,
집 밖 짝이 안 움직이면 ②③④도 침묵한다.)
\end{quote}

\section{챔피언은 그대로다}

\texttt{it} $220$ · lr $0.10$ · K $32$ 은 두 자 모두에서 최선이거나
비긴다. \textbf{용량 쪽은 다시 닫는다} --- 이번엔 자 둘로.

한 가지가 달라졌다. 노트 460 이 \texttt{it} $120$ 을 \emph{환율 $0.54$}
로 기각했는데, 무리 안 판으로 보면 $+0.0001$ 로 챔피언과 \textbf{비긴다}.
기각은 유지되지만(판을 $-0.0036$ 잃는다) \emph{잃는 것이 무리 사이
순위뿐}이라는 것을 이제 안다. 배포가 무리 안에서 고른다면 \texttt{it}
$120$ 은 공짜다.

\section*{남은 것}

\begin{itemize}
\item \textbf{깊이도 같은 무늬인가} --- 노트 472 에서 깊이 $12$ 는 판
$-0.0016$ · 무리 안 $-0.0023$ 으로 \emph{둘 다} 잃었다. \texttt{it} ·
lr 과 다른 자리다. 왜 깊이만 다른지 안 봤다.
\item \textbf{조건 ①$\sim$④ 에 무리 안 판을 어떻게 끼울까} --- 지금은
용량 변경에만 걸었다. 축 추가에도 걸어야 하나.
\item \textbf{배포 범위} --- 여전히 사람에게 물을 것. 이제 물을 이유가
하나 더 늘었다.
\end{itemize}

\end{document}
